% BEGIN R28_BOUNDARY_OCCUPANCY
% Insertar en MG01 tras la definición de c_j,o_j, antes de su valor.
La loi de la frontière de profondeur six fixe \(c_j=1\). L'axe \(z\) et l'
agrégat \(q\) étant déjà déterminés, les sommes entières latérales sont
\[
 (x_j+y_j)_{j=1}^6=(2,4,3,4,4,2).
\]
La règle d'occupation minimale sélectionne, parmi ses décompositions avec \(x_j,y_j\in\{0,1,2\}\), les colonnes ayant le plus petit nombre d'entrées non nulles, en comptant également \(z_j\). Pour une somme de deux, les paires extrêmes \((0,2),(2,0)\) occupent une position latérale ; pour une somme de trois, les paires \((1,2),(2,1)\) en occupent deux ; pour une somme de quatre, seule existe \((2,2)\). En ajoutant l'occupation de l'axe, on obtient \(o=(2,2,3,2,3,2)\). À cette frontière, avec \(s=[x+y]_3\), la même règle s'exprime comme
\[
 o_j=2+\mathbf1_{\{z_j\ne0,\ s_j\ne2\}}.
\]
% END R28_BOUNDARY_OCCUPANCY
